\section{Density Theorem}

Let A be a k-algebra. A left pseudomodule M is called semisimple if it is the direct sum of a family of simple left pseudomodules.

Lemma 1. --- Let M be a semisimple left A-pseudomodule. Let B be the bicommutant of the \(\widetilde{A}\) -module M. Then every A-sub-pseudomodule of M is a B-submodule of M.

The \(\widetilde{\mathrm{A}}\) -module M is semisimple. Let N be an A-sub-pseudomodule of M. Then N is an \(\widetilde{\mathrm{A}}\) -submodule of M. By Corollary 2 of VIII, p.~56, there exists a projector \(p\) of the A-module M with image N. Since we have the relation \(pb = bp\) for every \(b \in \mathrm{B}\) , we obtain that N is a B-submodule of M.

Lemma 2. --- Let M be a semisimple left A-pseudomodule, and let x be an element of M. There exists an element \(a \in A\) such that ax = x.

Let N be the left A-pseudomodule \(\widetilde{Ax}/Ax\) . It satisfies \(AN = \{0\}\) . By Corollary 3 of VIII, p.~56 applied to the \(\widetilde{A}\) -modules M and \(\widetilde{Ax}\) , the A-pseudomodule N is semisimple. By definition, every simple sub-pseudomodule S of N satisfies \(AS \neq \{0\}\) . Consequently, the pseudomodule N is zero, and we obtain that \(x \in Ax\) .

THEOREM 1 (Jacobson's density theorem). --- Let M be a semisimple left A-pseudomodule. Let b be an element of the bicommutant of the \(\widetilde{\mathrm{A}}\) -module M. Let \(\mathrm{F} = \{x_1, \ldots, x_n\}\) be a finite subset of M. Then there exists an element \(a \in \mathrm{A}\) such that \(bx_i = ax_i\) for every \(i \in [1, n]\) .

Let B be the bicommutant of the \(\widetilde{\mathrm{A}}\) -module M. The A-pseudomodule \(\mathrm{M}^n\) is semisimple. Let \(\boldsymbol{x} = (x_1, \ldots, x_n) \in \mathrm{M}^n\) . It follows from Lemma 2 that \(\boldsymbol{x} \in \mathrm{A}\boldsymbol{x}\) . The bicommutant of the \(\widetilde{\mathrm{A}}\) -module \(\mathrm{M}^n\) coincides with the homotheties of the B-module \(\mathrm{M}^n\) (VIII, p.~79, Proposition 2). By Lemma 1, the A-subpseudomodule \(\mathrm{A}\boldsymbol{x}\) of \(\mathrm{M}^n\) is therefore a B-submodule of \(\mathrm{M}^n\) . We therefore have the inclusion \(\mathrm{B}\boldsymbol{x} \subset \mathrm{A}\boldsymbol{x}\) . So there exists an \(a \in \mathrm{A}\) such that \(b\boldsymbol{x} = a\boldsymbol{x}\) . The result follows.
% semantic-fragments: legacy-input-seam
\relax{}
